\documentclass[12pt,a4paper]{ctexart}

% 页面设置
\usepackage[top=2cm, bottom=2cm, left=2.5cm, right=2cm]{geometry}
\usepackage{titlesec}
\usepackage{titletoc}
\usepackage{hyperref}
\usepackage{xcolor}
\usepackage{fancyhdr}
\usepackage{enumitem}
\usepackage{multicol}

% 页面样式
\pagestyle{fancy}
\fancyhf{}
\fancyhead[L]{\small 软件工程基础——详细目录}
\fancyhead[R]{\small 开卷考试参考}
\fancyfoot[C]{\thepage}
\renewcommand{\headrulewidth}{0.5pt}

% 标题格式
\titleformat{\section}{\Large\bfseries}{}{0em}{}
\titleformat{\subsection}{\large\bfseries}{}{0em}{}
\titleformat{\subsubsection}{\normalsize\bfseries}{}{0em}{}

\titlespacing{\section}{0pt}{8pt}{4pt}
\titlespacing{\subsection}{0pt}{6pt}{3pt}
\titlespacing{\subsubsection}{0pt}{4pt}{2pt}

% 紧凑列表
\setlist{nosep, leftmargin=2em, itemsep=0pt, parsep=0pt}
\setlength{\parskip}{0pt}
\setlength{\parindent}{0pt}

% 超链接设置
\hypersetup{
    colorlinks=true,
    linkcolor=black,
    urlcolor=blue
}

\begin{document}

\begin{center}
    {\Huge\bfseries 软件工程基础}\\[0.5cm]
    {\Large 详细考试目录}\\[0.3cm]
    {\normalsize （北京航空航天大学软件学院\quad 杨溢龙）}\\[0.5cm]
    \rule{\textwidth}{1pt}\\[0.5cm]
    {\normalsize 全册共186页\quad | \quad 8章内容\quad | \quad 按PDF页码标注}
\end{center}

\vspace{0.5cm}

% ============================================
% 第1章
% ============================================
\section{第1章：软件与软件工程 \hfill（第1--17页）}

\subsection{Part 1\quad 软件 \hfill（第1--7页）}

\subsubsection{1.1\quad 软件的基本概念与特点 \hfill（第1页）}
\begin{itemize}
    \item 软件的基本概念
    \item 软件的概念及特点
    \item 软件的分类
    \item 软件的地位和作用
    \item 建设狗窝和建设大厦的区别
    \item 算法 vs 程序 vs 软件
    \item 何为程序？——程序的两种形式（源代码/可执行代码）
\end{itemize}

\subsubsection{1.2\quad 程序编写与软件开发 \hfill（第2页）}
\begin{itemize}
    \item 程序如何编写？
    \item 软件开发难点与挑战
    \item 从程序到软件——工程开发的启发
    \item 目标领域类型（金融、信息技术、医疗等）
    \item 利益相关者（领域专家、设计师、架构师、开发/测试/运维工程师）
    \item 编程语言种类（C/C++/Java/Python/Dart/Go/Rust）
    \item 运行平台（Web/Desktop/Mobile）
    \item 系统类型（信息系统、嵌入式、CPS、区块链、AI使能系统）
\end{itemize}

\subsubsection{1.3\quad 软件的构成：程序+数据+文档 \hfill（第3页）}
\begin{itemize}
    \item 软件开发的再认识
    \item 软件开发核心问题（什么是对？什么是好？）
    \item 从需求到代码的分阶段跨越
    \item 何为软件？（程序 + 数据 + 文档）
    \item 何为数据？（初始化数据 + 测试数据 + 日常数据）
    \item 数据示例（用户、订单、交易、日志数据）
    \item 数据处理（表示、获取、存储、检索、分析）
\end{itemize}

\subsubsection{1.4\quad 软件文档 \hfill（第3--4页）}
\begin{itemize}
    \item 何为文档？（开发文档 + 管理文档 + 产品文档）
    \item 为什么需要文档？（记不住、理不清、讲不明）
    \item 编写文档目的（阐述清楚、发现问题、开展交流、促进管理）
    \item Agent工具生成能力对比（ATM/TakeOutMS/LoanPS案例）
    \item 集成测试成功率 vs 单元测试成功率 vs 错误迭代次数
    \item openspec / speckit / claude 对比数据
    \item 文档的重要性（企业成熟度的标志）
    \item 软件文档分类：开发文档、管理文档、产品文档
    \item 软件文档的覆盖关系（6层：目标程序→源程序→详细设计→概要设计→需求规格→用户需求→合同）
\end{itemize}

\subsubsection{1.5\quad 软件文档规范与模版 \hfill（第5页）}
\begin{itemize}
    \item 软件文档规范（国军标GB438B）
    \item 参考模版
\end{itemize}

\subsubsection{1.6\quad 软件的特点小结 \hfill（第5页）}
\begin{itemize}
    \item 逻辑性、设计开发、易变性、复杂性、缺陷的隐蔽性
\end{itemize}

\subsubsection{1.7\quad 软件质量 \hfill（第5--6页）}
\begin{itemize}
    \item 软件质量要素（1/3）：正确性、可靠性、健壮性、有效性
    \item 软件质量要素（2/3）：安全性、可维护性、可移植性、可重用性
    \item 软件质量要素（3/3）：可理解性、可信性、持续性、易用性、互操作性
\end{itemize}

\subsubsection{1.8\quad 软件分类 \hfill（第6--7页）}
\begin{itemize}
    \item 按功能分类：系统软件、支撑软件、应用软件
    \item 关键基础软件
    \item 开源软件（Open Source Software）
    \item 开源软件许可证：MIT、BSD、Apache、MPL、LGPL、GPL
\end{itemize}

\subsection{Part 2\quad 软件演化及危机 \hfill（第7--10页）}

\subsubsection{1.9\quad 软件发展与挑战 \hfill（第8页）}
\begin{itemize}
    \item 软件发展与挑战
    \item 传统的个体作坊式软件开发
    \item 个体作坊式创作带来的问题
    \item 软件开发面临的挑战：开发成本高、缺陷密度高、需求增长快、成功概率低、生产率低
\end{itemize}

\subsubsection{1.10\quad 软件危机 \hfill（第9--10页）}
\begin{itemize}
    \item 软件危机的出现
    \item 软件危机（Software Crisis）定义与表现
    \item 示例：IBM 360 OS软件开发
    \item 软件维护困难
    \item 质量难以保证
    \item 软件错误后果
    \item 软件灾难故事：ARIANE 5火箭爆炸
    \item 软件灾难故事：特斯拉Autopilot
    \item 软件灾难故事：波音737 MAX
    \item 大模型灾难故事
    \item 软件危机产生根源总结
    \item 解决危机的途径——软件工程
\end{itemize}

\subsection{Part 3\quad 软件工程 \hfill（第11--17页）}

\subsubsection{1.11\quad 软件工程基本概念 \hfill（第11--12页）}
\begin{itemize}
    \item 软件工程产生的动机
    \item 软件开发方式的改变
    \item 软件工程的定义（IEEE 1993）
    \item 软件工程的作用
\end{itemize}

\subsubsection{1.12\quad 软件工程的发展历程 \hfill（第12--14页）}
\begin{itemize}
    \item 20世纪50-60年代：萌芽阶段
    \item 20世纪70年代：结构化方法兴起
    \item 20世纪80年代：面向对象方法出现
    \item 20世纪90年代：UML统一建模语言、敏捷方法
    \item 21世纪前十年：敏捷普及、开源兴起
    \item 近十年（2011-至今）：DevOps、AI赋能
    \item 生成式软件开发 / 大模型
    \item 氛围编程（Vibe Coding）
    \item 提示词工程（Prompt Engineering）
    \item 智能体（Agent）
\end{itemize}

\subsubsection{1.13\quad 软件工程的知识体系 \hfill（第14页）}
\begin{itemize}
    \item IEEE SWEBOK V3.0 / V4.0（软件工程知识体系指南）
    \item 软件工程的多学科交叉
\end{itemize}

\subsubsection{1.14\quad 软件工程的目标与三要素 \hfill（第15页）}
\begin{itemize}
    \item 软件工程的目标
    \item 软件工程的三要素：过程（Process）、方法学（Methodology）、工具（Tool）
\end{itemize}

\subsubsection{1.15\quad 软件工程基本原则 \hfill（第16--17页）}
\begin{enumerate}
    \item 抽象和建模（Abstraction and Modeling）
    \item 模块化（Modularity）
    \item 软件重用（Software Reuse）
    \item 信息隐藏（Information Hiding）
    \item 关注点分离（Separation of Concerns）
    \item 分而治之（Divide and Conquer）
    \item 双向追踪原则（Bidirectional Traceability）
    \item 工具辅助（Tool Support）
\end{enumerate}

\subsubsection{1.16\quad 软件工程的基本观点与小结 \hfill（第17页）}

\vspace{0.3cm}
% ============================================
% 第2章
% ============================================
\section{第2章：软件过程与项目管理 \hfill（第18--29页）}

\subsection{Part 1\quad 软件过程概述 \hfill（第18--20页）}

\subsubsection{2.1\quad 软件过程与生命周期 \hfill（第18--19页）}
\begin{itemize}
    \item 软件过程（Software Process）定义
    \item 软件过程模型（Software Process Model）
    \item 软件生命周期模型概论
    \item 需求分析（Requirement Analysis）
    \item 概要设计（Architecture Design）
    \item 详细设计（Detailed Design）
    \item 编程实现（Implementation）
    \item 集成测试（Integration Test）
    \item 确认测试（Validation Test）
    \item 边写边改模型（Code and Fix）—— 最原始的模型
\end{itemize}

\subsection{Part 2\quad 软件过程模型 \hfill（第20--23页）}

\subsubsection{2.2\quad 瀑布模型 \hfill（第21页）}
\begin{itemize}
    \item 瀑布模型（Waterfall Model, 1970）
    \item 特点：阶段顺序执行、文档驱动
    \item 过程分解与精化的核心思想：增量 / 迭代
    \item 过程分解——功能维度
    \item 过程分解——时间维度
\end{itemize}

\subsubsection{2.3\quad 增量模型与快速原型 \hfill（第22页）}
\begin{itemize}
    \item 增量模型（Incremental Model）
    \item 快速原型模型（Rapid Prototyping）
    \item 螺旋模型（Spiral Model）—— Boehm 1988
\end{itemize}

\subsubsection{2.4\quad 敏捷方法与Scrum \hfill（第22--23页）}
\begin{itemize}
    \item 敏捷宣言（Agile Manifesto, 2001）—— 4条价值观
    \item Scrum过程
    \item Scrum角色：Product Owner（产品负责人）、Scrum Master、Team（开发团队）
    \item Sprint（冲刺）概念
\end{itemize}

\subsection{Part 3\quad 过程模型选择策略 \hfill（第23--24页）}
\begin{itemize}
    \item 软件过程模型的选择原则
    \item 场景一：V模型 / 外包项目
    \item 场景二：高风险项目 —— 螺旋模型
    \item 场景三：复刻产品 —— 增量模型
    \item 场景四：持续更新 —— 迭代模型 / 敏捷
    \item 模型选择小结
\end{itemize}

\subsection{Part 4\quad 软件项目管理 \hfill（第25--29页）}

\subsubsection{2.5\quad 软件项目概述 \hfill（第25--26页）}
\begin{itemize}
    \item 软件项目及其特点
    \item 软件项目的任务
    \item 软件项目实施方法——工程化
    \item 何为软件项目管理？（三要素：过程、人员、产品）
    \item 软件项目的管理要素
\end{itemize}

\subsubsection{2.6\quad 过程管理 \hfill（第27--28页）}
\begin{itemize}
    \item 过程定义
    \item 软件度量
    \item 成本预算
    \item 活动排序
    \item 软件项目计划
    \item 甘特图（Gantt Chart）—— 描述项目进度计划
    \item 项目跟踪
    \item 风险管理
\end{itemize}

\subsubsection{2.7\quad 产品管理 \hfill（第28--29页）}
\begin{itemize}
    \item 软件质量保证（SQA）
    \item 软件配置管理（SCM）
    \item 基线的概念（Baseline）
    \item 软件需求管理
\end{itemize}

\subsubsection{2.8\quad 人员管理 \hfill（第29页）}
\begin{itemize}
    \item 人员管理
    \item 软件项目团队
\end{itemize}

\vspace{0.3cm}
% ============================================
% 第3章 / 第5章 需求工程
% ============================================
\section{第3章 / 第5章：需求工程 \hfill（第30--79页）}

\subsection{Part 1\quad 需求基本概念 \hfill（第30--36页）}

\subsubsection{3.1\quad 需求的定义与重要性 \hfill（第30--32页）}
\begin{itemize}
    \item 软件开发的本质
    \item 需求重要性（价值所在 / 开发基础 / 验收标准）
    \item 空巢老人看护系统示例
    \item 软件需求定义（IEEE 1997标准）
    \item 何为需求？——解读IEEE标准
\end{itemize}

\subsubsection{3.2\quad 利益相关者分析 \hfill（第32--33页）}
\begin{itemize}
    \item 软件系统的利益相关者（用户、客户、系统分析员、开发者）
    \item 软件需求及其利益相关方
    \item 利益相关者对软件需求差异和冲突
\end{itemize}

\subsubsection{3.3\quad 软件需求的特点 \hfill（第34页）}
\begin{itemize}
    \item 隐式性（隐含的、未明确表达的）
    \item 隐晦性（难以明确表达的）
    \item 多源性（来源多样）
    \item 易变性（随时间变化）
    \item 领域知识相关性
    \item 价值不均性（不同需求价值不同）
\end{itemize}

\subsubsection{3.4\quad 软件需求的质量要求 \hfill（第34--35页）}
\begin{itemize}
    \item 有价值、正确、完整、无二义性
    \item 可行、一致、可追踪、可验证
\end{itemize}

\subsubsection{3.5\quad 可行性分析 \hfill（第35--36页）}
\begin{itemize}
    \item 技术可行性
    \item 设备可行性
    \item 进度可行性
    \item 成本可行性
    \item 商业可行性
    \item 社会可行性（法律法规）
    \item 可行性研究报告
\end{itemize}

\subsection{Part 2\quad 需求工程 \hfill（第36--41页）}

\subsubsection{3.6\quad 需求工程的定义与特点 \hfill（第36--37页）}
\begin{itemize}
    \item 需求工程的重要性
    \item 何为需求工程？（Requirements Engineering）
    \item 需求工程的特点
    \item ISO/IEC/IEEE 15288 / 12207 / 29148 标准
    \item 需求工程指南及参考书目（IREB、INCOSE、IEEE 830）
\end{itemize}

\subsubsection{3.7\quad 需求工程过程模型 \hfill（第38--40页）}
\begin{itemize}
    \item 螺旋模型 —— 三层次三阶段
    \item 需求的三个层次：业务需求、用户需求、系统需求
    \item 系统需求：功能性需求 vs 非功能性需求
    \item 非功能性需求 —— ISO 25010标准分类
    \item 需求工程三阶段：需求获取、需求表示/精化与规约、需求确认与验证
\end{itemize}

\subsubsection{3.8\quad 需求管理 \hfill（第40--41页）}
\begin{itemize}
    \item 软件需求的缺陷
    \item 软件需求变更管理
    \item 需求工程相关的CASE工具
\end{itemize}

\subsection{Part 3\quad 需求获取 \hfill（第41--50页）}

\subsubsection{3.9\quad 需求获取的来源与方法 \hfill（第41--43页）}
\begin{itemize}
    \item 软件需求从何而来？
    \item 需求获取为什么难？
    \item 成立需求分析的联合工作小组（JRT）
    \item 需求获取的方法：
    \begin{itemize}
        \item 访谈和会议
        \item 调查问卷
        \item 现场观摩
        \item 分析业务资料
        \item 群体化方法（Brainstorming / JAD）
        \item 大语言模型推荐
    \end{itemize}
\end{itemize}

\subsubsection{3.10\quad 需求获取的一般过程（1/4）——明确问题与边界 \hfill（第44页）}
\begin{itemize}
    \item 明确软件要解决的问题
    \item 课程实践创意示例：无人值守图书馆
    \item 课程实践创意示例：多无人机联合搜寻
\end{itemize}

\subsubsection{3.11\quad 需求获取的一般过程（2/4）——识别利益相关方与功能需求 \hfill（第45--46页）}
\begin{itemize}
    \item 识别利益相关方
    \item 获取软件的功能性需求
    \item 合理的软件需求
    \item 获取和构思非功能性需求
\end{itemize}

\subsubsection{3.12\quad 需求获取的一般过程（3/4）——描述初步软件需求 \hfill（第46--50页）}
\begin{itemize}
    \item 自然语言描述及不足
    \item 结构化自然语言描述
    \item 用户故事（User Story）—— 三要素：角色、功能、价值
    \item EARS（Easy Approach to Requirements Syntax）—— 简洁需求句法
    \item EARS五类模板：Ubiquitous / Event-driven / State-driven / Unwanted / Optional
    \item 交互流程描述（场景描述）
    \item 基本流程（Basic Flow）
    \item 扩展流程（Extension Flow）——编号编写、异常处理、选择
    \item 撰写初步软件需求描述文档
\end{itemize}

\subsubsection{3.13\quad 需求获取的一般过程（4/4）——评审初步需求 \hfill（第50页）}
\begin{itemize}
    \item 评审初步软件需求
    \item 需求获取（总图）
\end{itemize}

\subsection{Part 4\quad 需求精化与规约 \hfill（第51--76页）}

\subsubsection{3.14\quad 需求精化与优先级 \hfill（第51--52页）}
\begin{itemize}
    \item 为什么要精化初始软件需求
    \item 软件需求的精化与规约的任务
    \item 分析和确定软件需求优先级
    \item 分析软件需求重要性
    \item 示例：空巢老人看护软件的需求重要性分析
    \item 确定需求迭代先后次序
\end{itemize}

\subsubsection{3.15\quad 需求精化的核心问题 \hfill（第53页）}
\begin{itemize}
    \item 两层次跨越：从问题域到解域
\end{itemize}

\subsubsection{3.16\quad 面向对象需求建模 \hfill（第54--55页）}
\begin{itemize}
    \item 面向对象系统需求的建模方法
    \item 用例模型（行为视角）
    \item 领域概念模型（结构视角）
    \item 系统操作合约
    \item 一阶逻辑 / 对象约束语言（OCL）
\end{itemize}

\subsubsection{3.17\quad UML统一建模语言 \hfill（第56--57页）}
\begin{itemize}
    \item UML: Unified Modeling Language
    \item Unified（统一）、Modeling（建模）、Language（语言）
    \item UML本质 —— 是一种语言，不是方法学
    \item UML与软件开发方法、过程的关系
    \item 面向对象建模语言
    \item 面向对象需求分析方法的优势和特色
    \item UML应用范围
    \item UML图（9种图 / 5类）
    \item UML的视点及图：
    \begin{itemize}
        \item 用例图（Use Case Diagram）
        \item 静态图：类图、对象图、包图
        \item 结构：构件图、部署图
        \item 行为：状态图、活动图、通信图、顺序图
    \end{itemize}
\end{itemize}

\subsubsection{3.18\quad 用例图（Use Case Diagram）\hfill（第57--61页）}
\begin{itemize}
    \item 参与者（Actor）—— 识别参与者
    \item 参与者之间的泛化关系
    \item 系统边界框
    \item 用例（Use Case）
    \item 用例的特征 / 特点
    \item 用例往往是动宾短语
    \item 用例独立和完备性
    \item 用例由参与者启动
    \item 用例要有可观测的执行结果
    \item 用例描述
    \item 用例的识别
    \item 执行者与用例间的关系
    \item 用例间的关系：
    \begin{itemize}
        \item 包含（Include）—— 必须执行
        \item 扩展（Extend）—— 可选执行
    \end{itemize}
    \item 超市管理系统——用例模型示例
\end{itemize}

\subsubsection{3.19\quad 顺序图（Sequence Diagram）\hfill（第62--64页）}
\begin{itemize}
    \item 顺序图的概念
    \item 三种主要作用
    \item 顺序图的组成元素：对象、生命线、激活、消息
    \item 对象（Object）
    \item 生命线（Lifeline）
    \item 激活（Execution / Activation）
    \item 消息（Message）
    \item 同步消息（Synchronous） vs 异步消息（Asynchronous）
    \item 系统顺序图（需求） VS 对象顺序图（设计）—— 重要区分
\end{itemize}

\subsubsection{3.20\quad 面向对象的核心概念 \hfill（第64--66页）}
\begin{itemize}
    \item 对象（Object）
    \item 类（Class）—— 对象与类的关系
    \item 继承（Inheritance）
    \item 多态（Polymorphism）
    \item 覆盖/重写（Override）
    \item 重载（Overload）
    \item 消息（Message）
    \item 聚合（Aggregation）和组合（Composition）—— 整体-部分关系
\end{itemize}

\subsubsection{3.21\quad 领域概念模型 \hfill（第66--67页）}
\begin{itemize}
    \item 识别类与对象 —— 系统实体的识别方法
    \item 识别属性
    \item 识别关联关系
    \item 关联名称 / 角色 / 导向性
    \item 关联类
    \item 关联的多重性（1:1, 1:n, n:m）
    \item 领域模型（Domain Model）
    \item 小结
\end{itemize}

\subsubsection{3.22\quad 系统操作合约 \hfill（第68--69页）}
\begin{itemize}
    \item 系统操作（System Operation）
    \item 系统操作合约（System Operation Contract）
    \item 操作签名、交叉引用、前置条件、后置条件
    \item 面向对象的基本操作（创建对象、连接变化、属性变化）
    \item OCL（Object Constraint Language）—— 对象约束语言
    \item OCL合约
\end{itemize}

\subsubsection{3.23\quad 面向过程的需求建模（结构化分析）\hfill（第70--75页）}
\begin{itemize}
    \item 结构化分析思想
    \item 数据流模型（行为视角）
    \item 实体关系模型（结构视角）
    \item 结构化分析方法概述
    \item 指导性原则
    \item 结构化方法的局限性
    \item 结构化分析工具：DFD、DD、结构化语言、判定表、判定树
\end{itemize}

\subsubsection{3.24\quad 数据流图（DFD）\hfill（第71--72页）}
\begin{itemize}
    \item 数据流图（Data Flow Diagram）
    \item 基本成份：外部实体、处理、数据流、数据存储
    \item 环境图（0级DFD / 顶层图）
    \item 数据流图的绘制步骤
    \item 分层数据流图
    \item 分层原则
    \item 特征与用途
    \item 教学管理系统分层数据流图示例
\end{itemize}

\subsubsection{3.25\quad 数据字典 \hfill（第73页）}
\begin{itemize}
    \item 数据字典（Data Dictionary）的定义与用途
    \item 数据字典的种类
    \item 数据字典的定义方法
    \item 数据字典——数据流
    \item 数据字典——数据项
    \item 数据字典——数据存储
\end{itemize}

\subsubsection{3.26\quad 结构化语言与判定工具 \hfill（第73--74页）}
\begin{itemize}
    \item 结构化语言
    \item 判定表（Decision Table）
    \item 判定树（Decision Tree）
\end{itemize}

\subsubsection{3.27\quad 图书借阅管理系统——完整DFD示例 \hfill（第74--75页）}
\begin{itemize}
    \item 顶层数据流图
    \item 中层数据流图
    \item 各模块底图（处理数据流图）
\end{itemize}

\subsubsection{3.28\quad OO vs 面向过程建模对比 \hfill（第76页）}
\begin{itemize}
    \item 基于模型的系统需求建模——两种方法的比较
\end{itemize}

\subsection{Part 5\quad 需求确认与验证 \hfill（第76--79页）}

\subsubsection{3.29\quad 需求验证 \hfill（第76--77页）}
\begin{itemize}
    \item 解决软件需求问题（遗漏、无源头、不一致、不正确、不规范）
    \item 需求验证目标
    \item Validation（确认） vs Verification（验证）—— 重要区分
    \item 需求确认与验证的手段：
    \begin{itemize}
        \item 需求评审（Review）
        \item 原型生成（Prototyping）
        \item 测试用例生成（Test Case Generation）
    \end{itemize}
    \item 软件需求评审的步骤
    \item 评审软件需求
\end{itemize}

\subsubsection{3.30\quad 需求规格说明 \hfill（第78--79页）}
\begin{itemize}
    \item 软件需求分析的输出：需求模型 + 软件原型 + 需求文档
    \item 需求规格说明（SRS）
    \item 用户需求报告
    \item 需求分析规格说明书
    \item 撰写软件需求文档的注意事项
    \item 软件需求文档模板
\end{itemize}

\subsubsection{3.31\quad 需求管理活动 \hfill（第79页）}
\begin{itemize}
    \item 需求工程管理活动（变更控制、版本控制、需求状态跟踪、需求追踪）
\end{itemize}

\vspace{0.3cm}
% ============================================
% 第4章
% ============================================
\section{第4章：软件概要设计 \hfill（第80--107页）}

\subsection{Part 1\quad 软件设计概述 \hfill（第80--85页）}

\subsubsection{4.1\quad 软件设计的基本概念 \hfill（第80--82页）}
\begin{itemize}
    \item 从需求到设计和编码
    \item 软件设计的基本概念
    \item 何为软件设计？（定义、内容）
    \item 软件设计的内容
    \item 软件设计特点：多样性和差异性
    \item 需求分析、软件设计、软件实现间的关系
    \item 软件设计的质量要求：正确性、充分性、优化性、简单性
\end{itemize}

\subsubsection{4.2\quad 软件设计基本原则 \hfill（第83--85页）}
\begin{enumerate}
    \item \textbf{抽象原则}（Abstraction）—— 抽象在软件设计中的应用
    \begin{itemize}
        \item 软件设计抽象层次的变化
        \item 体系结构层次 / 构件层次的设计抽象
    \end{itemize}
    \item \textbf{模块化}（Modularity）—— 模块分解与开发成本的关系
    \item \textbf{高内聚度原则}（High Cohesion）
    \begin{itemize}
        \item 内聚度分类（7种）：偶然内聚、逻辑内聚、时间内聚、过程内聚、通讯内聚、顺序内聚、功能内聚
    \end{itemize}
    \item \textbf{低耦合度原则}（Low Coupling）
    \begin{itemize}
        \item 耦合度分类（7种）：内容耦合、公共耦合、外部耦合、控制耦合、标记耦合、数据耦合、无直接耦合
    \end{itemize}
    \item \textbf{信息隐藏原则}（Information Hiding）—— 示例
    \item \textbf{关注点分离原则}（Separation of Concerns）
    \item \textbf{软件重用原则}（Software Reuse）
\end{enumerate}

\subsection{Part 2\quad 软件设计过程 \hfill（第85--89页）}

\subsubsection{4.3\quad 软件设计的过程 \hfill（第85--87页）}
\begin{itemize}
    \item 软件设计过程概述
    \item 软件设计元素：子系统、软件组件、设计类/设计函数
    \item 软件设计的过程——体系结构设计
    \item 分层体系结构风格示例
    \item 软件设计的过程——详细设计
\end{itemize}

\subsubsection{4.4\quad 软件设计的输出与评审 \hfill（第88--89页）}
\begin{itemize}
    \item 软件设计的输出：设计模型 + 设计文档
    \item 软件设计文档的规范
    \item 软件设计中的缺陷类型
    \item 软件设计的评审
    \item 软件设计的CASE工具
\end{itemize}

\subsection{Part 3\quad 软件架构设计 \hfill（第89--94页）}

\subsubsection{4.5\quad 软件体系结构设计概述 \hfill（第89--90页）}
\begin{itemize}
    \item 软件架构设计概述
    \item 软件体系结构设计示意图
    \item 软件体系结构设计的目标
    \item 软件体系结构设计的任务
    \item 与软件需求的关系
    \item 与详细设计的关系
    \item 软件体系结构设计的过程（4步骤）
\end{itemize}

\subsubsection{4.6\quad 设计初步的软件体系结构 \hfill（第91--92页）}
\begin{itemize}
    \item 设计初步的软件体系结构
    \item 不同体系结构模式适合的应用类型
    \item 重用开源软件及已有软件资产
    \item 到哪里找开源软件？（GitHub / SourceForge / Gitee）
    \item 示例：搜寻和重用软件资产
\end{itemize}

\subsubsection{4.7\quad 精化软件体系结构 \hfill（第93--94页）}
\begin{itemize}
    \item 精化软件体系结构
    \item 确定子系统及其接口
    \item 确定构件及其接口
    \item 优化体系结构设计
    \item 撰写软件体系结构设计文档
    \item 评审软件体系结构设计
\end{itemize}

\subsection{Part 4\quad 架构设计模型 \hfill（第95--101页）}

\subsubsection{4.8\quad 面向对象的架构模型 \hfill（第95--98页）}
\begin{itemize}
    \item 回顾：面向对象的需求模型
    \item 软件体系结构的概念
    \item 面向对象软件设计过程
    \item \textbf{包图（Package Diagram）}
    \begin{itemize}
        \item 包图的基本概念
        \item 包的作用
        \item 包的依赖关系
        \item 示例：包图表示体系结构
    \end{itemize}
    \item \textbf{构件图（Component Diagram）}
    \begin{itemize}
        \item 组件的概念
        \item 组件及其接口
        \item 组件和接口的表示（供给接口 / 需求接口 / 端口）
        \item 示例：软件组件（.dll / .jar）
        \item 软件体系结构的逻辑视图：组件图
    \end{itemize}
\end{itemize}

\subsubsection{4.9\quad 面向过程的架构模型 \hfill（第99--101页）}
\begin{itemize}
    \item 面向过程概要设计的主要内容
    \item H图（层次图 / Hierarchy Chart）
    \item IPO图（输入-处理-输出图）
    \item 结构图（Structure Chart）
    \item 面向数据流的设计方法
    \item 变换分析（Transform Analysis）
    \item 事物分析（Transaction Analysis）
\end{itemize}

\subsection{Part 5\quad 软件架构设计模式 \hfill（第101--107页）}

\subsubsection{4.10\quad 软件体系结构风格 \hfill（第101--102页）}
\begin{itemize}
    \item 软件体系结构风格（Architecture Style）
    \item 建筑风格类比（中式/西式）
    \item 软件设计模式（Design Pattern）
    \item 常用软件体系结构风格概述
\end{itemize}

\subsubsection{4.11\quad 分层体系结构模式 \hfill（第102--105页）}
\begin{itemize}
    \item 分层体系结构模式（Layered Architecture）
    \item 分层模式的约束 / 特点
    \item 示例：网络OSI模型 / 云计算架构
    \item 三层结构设计（Three-Tier Architecture）
    \begin{itemize}
        \item 表示层（Presentation Layer）
        \item 中间层 / 业务逻辑层（Business Logic Layer）
        \item 数据层（Data Layer）
    \end{itemize}
    \item 三层协调工作 / 三层结构优点
    \item \textbf{MVC风格}（Model-View-Controller）
    \begin{itemize}
        \item Model（模型）、View（视图）、Controller（控制器）
        \item MVC风格的约束
    \end{itemize}
    \item \textbf{MVP模式}（Model-View-Presenter）
    \begin{itemize}
        \item MVP与MVC的区别
    \end{itemize}
    \item \textbf{MVVM框架}（Model-View-ViewModel）
    \item 贫血模型（Anemic Model）与充血模型（Rich Model）
\end{itemize}

\subsubsection{4.12\quad 管道与过滤器风格 \hfill（第106页）}
\begin{itemize}
    \item 管道与过滤器风格（Pipe-and-Filter）
    \item 模式的特点 / 约束
    \item 示例：Unix Shell管道
\end{itemize}

\subsubsection{4.13\quad SOA与消息总线 \hfill（第106--107页）}
\begin{itemize}
    \item 面向服务的体系结构（SOA）
    \item SOA风格的特点
    \item 消息总线风格（Message Bus）
    \item ESB（Enterprise Service Bus）体系结构
\end{itemize}

\vspace{0.3cm}
% ============================================
% 第5章
% ============================================
\section{第5章：软件详细设计 \hfill（第108--134页）}

\subsection{Part 1\quad 详细设计概述 \hfill（第108--111页）}

\subsubsection{5.1\quad 详细设计的任务与定位 \hfill（第108--110页）}
\begin{itemize}
    \item 面向对象/面向过程方法的软件设计过程（回顾）
    \item 详细设计的动机
    \item 详细设计的任务与定位
    \item 目标与原则
    \item 主要活动
    \item \textbf{概要设计 vs 详细设计}（重要对比）：
    \begin{itemize}
        \item 概要设计：做什么（What）—— 体系结构层面
        \item 详细设计：怎么做（How）—— 实现细节层面
    \end{itemize}
    \item 软件详细设计的要求和原则
\end{itemize}

\subsubsection{5.2\quad 软件详细设计的活动 \hfill（第110--111页）}
\begin{enumerate}
    \item 用例设计
    \item 子系统/软构件设计
    \item 类/函数设计
    \item 持久化设计（数据设计）
\end{enumerate}

\subsection{Part 2\quad 面向对象的详细设计过程 \hfill（第111--126页）}

\subsubsection{5.3\quad 详细设计的过程（4步骤）\hfill（第112--113页）}
\begin{itemize}
    \item 步骤1：用例设计
    \item 用例设计的任务
    \item 系统与设计类间对应关系
    \item 设计元素及其交互
    \item 示例：用例顺序图 / 设计方案
\end{itemize}

\subsubsection{5.4\quad 类设计 \hfill（第114--120页）}
\begin{itemize}
    \item 类设计的任务
    \item 类设计原则：准确化、细节化、一致性、遵循设计原则
    \item 注意事项
    \item 类设计过程：
    \begin{enumerate}
        \item 确定类的可见范围（public / protected / private）
        \item 精化类间的关系（继承 / 组合 / 聚合 / 关联 / 依赖）
        \item 精化类属性的设计
        \item 精化类方法的设计
        \item 构造类对象的状态图 / 活动图
        \item 评审和优化
    \end{enumerate}
    \item \textbf{活动图（Activity Diagram）}
    \begin{itemize}
        \item 基本概念
        \item 组成元素：动作和活动节点、开始和终止、控制流、判断节点、分叉节点、合并节点、结合节点
        \item 对象流（Object Flow）
        \item 泳道（Swimlane）
    \end{itemize}
    \item \textbf{状态图（State Machine Diagram）}
    \begin{itemize}
        \item 概念、主要图符
        \item 转移（Transition）：触发器→守卫→影响
    \end{itemize}
    \item 类设计输出的软件制品
\end{itemize}

\subsubsection{5.5\quad 持久化设计（数据设计）\hfill（第121--125页）}
\begin{itemize}
    \item 为什么要进行持久化设计
    \item 数据设计的任务
    \item 数据设计的原则
    \item 持久化设计的任务
    \item 数据设计过程：
    \begin{enumerate}
        \item 确定永久数据（哪些数据需要持久化存储）
        \item 确定持久数据的存储和组织方式
        \item 确定持久数据条目
        \item 关系映射：
        \begin{itemize}
            \item 1:n关联关系的映射（外键）
            \item n:m关联关系的映射（中间表）
            \item 继承关系的数据库表设计（3种策略）
        \end{itemize}
        \item 设计数据操作（CRUD）
        \item 评审和优化数据设计
    \end{enumerate}
    \item 示例：设计持久数据
    \item 示例：设计永久数据的操作
\end{itemize}

\subsubsection{5.6\quad 设计整合与验证 \hfill（第125--126页）}
\begin{itemize}
    \item 软件详细设计的输出
    \item 设计整合
    \item 设计验证
\end{itemize}

\subsection{Part 3\quad 面向过程的详细设计 \hfill（第126--133页）}

\subsubsection{5.7\quad 过程设计工具 \hfill（第126--130页）}
\begin{itemize}
    \item 面向过程详细设计概述
    \item \textbf{程序流程图（Flow Chart）}
    \begin{itemize}
        \item 常见的程序流程图结构：顺序、选择、循环
    \end{itemize}
    \item \textbf{N-S图（Nassi-Schneiderman Diagram / 盒图）}
    \begin{itemize}
        \item 特点：结构化、无GOTO
    \end{itemize}
    \item \textbf{PDL（Program Design Language / 程序设计语言）}
    \begin{itemize}
        \item 伪代码形式
    \end{itemize}
    \item \textbf{PAD（Problem Analysis Diagram / 问题分析图）}
    \item \textbf{决策表（Decision Table）}
    \item \textbf{判定树（Decision Tree）}
\end{itemize}

\subsubsection{5.8\quad Jackson方法（面向数据结构的设计方法）\hfill（第130--133页）}
\begin{itemize}
    \item Jackson方法的基本思想
    \item Jackson方法的3种基本类型：顺序、选择、重复
    \item Jackson图（树形结构）
    \item 改进的Jackson图
    \item Jackson方法解决问题的例子（零件库房管理）
    \item Jackson方法步骤（1-5步）
    \item 伪代码表示程序
\end{itemize}

\subsection{Part 4\quad 设计文档与评审 \hfill（第133--134页）}

\subsubsection{5.9\quad 撰写与评审设计文档 \hfill（第133--134页）}
\begin{itemize}
    \item 撰写设计文档（目录结构）
    \item 软件设计方案提纲
    \item 优化和评审用例设计方案
    \item 设计评审人员（用户代表、设计人员、程序员、需求人员、质量人员、测试人员、配置管理员）
    \item 评审设计文档的要点：规范性、可追踪性、一致性、正确性、高质量、可实施性
    \item 评审步骤
    \item 小结
\end{itemize}

\vspace{0.3cm}
% ============================================
% 第6章
% ============================================
\section{第6章：软件实现 \hfill（第135--152页）}

\subsection{Part 1\quad 软件实现概论 \hfill（第135--140页）}

\subsubsection{6.1\quad 软件实现的概念与任务 \hfill（第135--136页）}
\begin{itemize}
    \item 何为软件实现（微观概念 / 与设计和测试的关系）
    \item 软件实现的任务
    \item 软件实现与软件设计之间的关系
    \item 软件实现的特点：生产性活动 + 创作性活动
\end{itemize}

\subsubsection{6.2\quad 编码原则与规范 \hfill（第137--138页）}
\begin{itemize}
    \item 编写代码的原则（1/3）：易读、易改
    \item 编写代码的原则（2/3）：异常处理、代码与模型一致
    \item 编写代码的原则（3/3）：降低复杂度、软件重用
    \item 编程规范（Style Guide）
    \item 遵循编码风格（1/4--4/4）：格式化、简洁性、文档、结构化
    \item \textbf{代码重构（Code Refactoring）}—— 改善内部结构而不改变外部行为
\end{itemize}

\subsubsection{6.3\quad 软件实现的过程与工具 \hfill（第139--140页）}
\begin{itemize}
    \item 软件实现的过程：编码 + 单元测试 + 调试
    \item 支持软件实现的CASE工具 —— IDE
    \item 目前主流IDE及其特点：VSCode、JetBrains系列、Eclipse
    \item 目前主流IDE发展历史
    \item 未来IDE特点与趋势：服务化、云端化、定制化
\end{itemize}

\subsection{Part 2\quad 编程语言 \hfill（第140--143页）}

\subsubsection{6.4\quad 程序设计语言的类别 \hfill（第140--141页）}
\begin{itemize}
    \item 机器语言（Machine Language）
    \item 汇编语言（Assembly Language）
    \item 结构化程序设计语言（C、Pascal等）
    \item 面向对象程序设计语言（Java、C++等）
    \item 描述性程序设计语言（SQL、HTML等）
    \item 程序设计语言的表达能力分类：命令式 / 函数式 / 逻辑式
\end{itemize}

\subsubsection{6.5\quad 主流编程语言 \hfill（第141--142页）}
\begin{itemize}
    \item C / C++语言特点
    \item Java语言特点
    \item Python语言特点
    \item Swift / Kotlin / ArkTS / TypeScript / Rust / 仓颉
    \item 不同编程语言Top 10排名
\end{itemize}

\subsubsection{6.6\quad 程序设计语言的选择 \hfill（第143页）}
\begin{itemize}
    \item 软件的应用领域
    \item 与遗留软件系统的交互
    \item 软件的特殊功能及需求
    \item 软件的目标平台
    \item 程序员的编程经验
\end{itemize}

\subsection{Part 3\quad 软件编码过程 \hfill（第143--149页）}

\subsubsection{6.7\quad 编写代码的任务 \hfill（第144--145页）}
\begin{enumerate}
    \item 编写实现类的代码
    \item 编写实现类方法的代码
    \item 编写实现类间关联的代码
    \item 编写实现聚合和组合关系的代码
    \item 编写实现接口关系的代码
    \item 编写实现继承关系的程序代码
    \item 编写实现包的代码
\end{enumerate}

\subsubsection{6.8\quad 界面与数据编码 \hfill（第146--147页）}
\begin{itemize}
    \item 编写用户界面代码
    \item 示例：LoginUI的编码实现
    \item 编写数据设计代码
    \item 示例：T\_User表的创建
    \item 示例：连接和关闭数据库的代码
    \item 示例：操作数据库的程序代码
\end{itemize}

\subsubsection{6.9\quad 框架与案例 \hfill（第148--149页）}
\begin{itemize}
    \item 什么是框架？（Framework）
    \item 框架 vs 类库的区别
    \item Case Study：基于Java技术栈的系统实现
    \item 例1：基于Java EE的单体式应用实现（分层模型）
    \item 设计模型到Java EE代码的映射过程
    \item 类图 / 顺序图映射
\end{itemize}

\subsubsection{6.10\quad 微服务架构补充 \hfill（第151--152页）}
\begin{itemize}
    \item 微服务架构模型
    \item 基于SpringCloud的微服务应用
    \item 微服务组件图 / 概念类图 / 顺序图
    \item Spring Cloud架构
    \item 服务发现配置
    \item 网关实现
    \item Spring Cloud应用在阿里云部署
\end{itemize}

\subsection{Part 4\quad 软件文档 \hfill（第149--150页）}
\begin{itemize}
    \item 软件实现文档
    \item 用户指南的编写要求
    \item 用户指南的排版建议
    \item 用户指南的编写原则
    \item 软件实现管理文档
    \item 小结
\end{itemize}

\vspace{0.3cm}
% ============================================
% 第7章
% ============================================
\section{第7章 / 第6章续：软件测试 \hfill（第153--175页）}

\subsection{Part 1\quad 软件缺陷 \hfill（第153--156页）}

\subsubsection{7.1\quad 软件缺陷基本概念 \hfill（第153--154页）}
\begin{itemize}
    \item 故障（Fault）分类：按发生阶段、系统边界、现象、目标、目的、持久性
    \item 软件缺陷（Software Defect / Bug）
    \item 故障（Fault） → 错误（Error） → 失效（Failure）—— 缺陷传播链
    \item 软件缺陷的危害（经济损失、人身安全等）
\end{itemize}

\subsubsection{7.2\quad 缺陷描述与管理 \hfill（第155--156页）}
\begin{itemize}
    \item 软件缺陷不可避免的原因
    \item 软件缺陷的描述：标识符、类型、严重程度、症状、修复优先级、状态、发现者、源头、原因
    \item 软件缺陷的应对方法：发现缺陷、修复缺陷、容忍缺陷、预防缺陷
    \item 程序调试（Debugging）—— 定义与步骤
    \item 缺陷管理工具（PingCode / Jira / ClickUp等）
\end{itemize}

\subsection{Part 2\quad 软件测试概论 \hfill（第157--160页）}

\subsubsection{7.3\quad 测试的基本概念 \hfill（第157--158页）}
\begin{itemize}
    \item 软件测试概论
    \item 在程序代码中找出软件缺陷
    \item 程序构成及其缺陷的潜在位置
    \item 软件测试的目的（发现错误、验证需求、提高质量）
    \item 软件测试的功效（不能证明程序无错）
    \item 软件测试的输入与输出
\end{itemize}

\subsubsection{7.4\quad 测试的本质与测试用例 \hfill（第159页）}
\begin{itemize}
    \item 软件测试的本质
    \item \textbf{测试用例（Test Case）}—— 四元组：
    \begin{enumerate}
        \item 输入数据
        \item 前置条件
        \item 操作步骤
        \item 预期输出
    \end{enumerate}
    \item 示例：测试用例的设计（"用户登录"）
    \item 软件测试的步骤（4步）
\end{itemize}

\subsubsection{7.5\quad 软件测试原则 \hfill（第160页）}
\begin{enumerate}
    \item 测试应有计划、有组织地进行
    \item 测试应追溯需求
    \item Pareto原则（80\%的错误来自20\%的模块）
    \item 穷举测试是不可能的
    \item 第三方测试更有效
    \item 测试用例应长期保留（回归测试用）
    \item 避免测试自己编写的程序
    \item 不仅要测试正常情况，还要测试异常情况
\end{enumerate}

\subsection{Part 3\quad 软件测试分类 \hfill（第160--166页）}

\subsubsection{7.6\quad 测试分类概述 \hfill（第161页）}
\begin{itemize}
    \item 静态测试 vs 动态测试
    \item 白盒测试 vs 黑盒测试 vs 灰盒测试
    \begin{itemize}
        \item 白盒测试：知道内部结构，测试内部逻辑
        \item 黑盒测试：不知道内部结构，测试功能
        \item 灰盒测试：部分知道内部结构
    \end{itemize}
\end{itemize}

\subsubsection{7.7\quad 白盒测试（White-Box Testing）\hfill（第161--164页）}
\begin{itemize}
    \item 白盒测试用例设计的指导原则
    \item \textbf{语句覆盖（Statement Coverage）}
    \begin{itemize}
        \item 每条语句至少执行一次
    \end{itemize}
    \item \textbf{判定覆盖（Branch Coverage / 分支覆盖）}
    \begin{itemize}
        \item 每个判定的每个分支至少执行一次
    \end{itemize}
    \item \textbf{组合覆盖（Condition Coverage）}
    \begin{itemize}
        \item 每个条件的所有可能组合至少执行一次
    \end{itemize}
    \item \textbf{基本路径测试（Basis Path Testing）}
    \begin{itemize}
        \item 步骤1：根据程序逻辑画出流程图
        \item 步骤2：将流程图转换为流图（Flow Graph）
        \begin{itemize}
            \item 流图中的判定点不应含复合条件
        \end{itemize}
        \item 步骤3：确定基本路径集合
        \begin{itemize}
            \item 圈复杂度 V(G) = E - N + 2（或 P + 1，或 判定节点数 + 1）
        \end{itemize}
        \item 步骤4：针对测试路径设计测试用例
        \item 步骤5：运行程序检验测试用例
    \end{itemize}
\end{itemize}

\subsubsection{7.8\quad 黑盒测试（Black-Box Testing）\hfill（第165--166页）}
\begin{itemize}
    \item 黑盒测试的思想、依据、特点
    \item \textbf{等价分类法（Equivalence Partitioning）}
    \begin{itemize}
        \item 有效等价类
        \item 无效等价类
        \item 实例：等价类划分法（三角形问题）
    \end{itemize}
    \item \textbf{边界值分析法（Boundary Value Analysis）}
    \begin{itemize}
        \item 原则：边界值、边界值±1
        \item 示例
    \end{itemize}
\end{itemize}

\subsection{Part 4\quad 软件测试过程 \hfill（第167--172页）}

\subsubsection{7.9\quad 软件测试活动层次 \hfill（第167页）}
\begin{itemize}
    \item 软件测试活动（5个层次）：
    \begin{enumerate}
        \item 单元测试（Unit Testing）
        \item 集成测试（Integration Testing）
        \item 确认测试（Validation Testing）
        \item 系统测试（System Testing）
        \item 验收测试（Acceptance Testing）
    \end{enumerate}
\end{itemize}

\subsubsection{7.10\quad 单元测试 \hfill（第167--168页）}
\begin{itemize}
    \item 程序单元测试的依据和对象
    \item 谁负责开展程序单元测试？（通常由开发人员）
    \item 单元测试技术（驱动模块 / 桩模块）
    \item 单元测试工具（JUnit / C++Test / pytest等）
\end{itemize}

\subsubsection{7.11\quad 集成测试 \hfill（第168--169页）}
\begin{itemize}
    \item 集成测试的对象、技术、内容
    \item \textbf{集成测试策略：}
    \begin{itemize}
        \item 大爆炸集成（Big Bang）—— 一次性集成所有模块
        \item 自顶向下集成（Top-Down）—— 需要桩模块
        \item 自底向上集成（Bottom-Up）—— 需要驱动模块
        \item 三明治集成（Sandwich）—— 结合自顶向下和自底向上
    \end{itemize}
\end{itemize}

\subsubsection{7.12\quad 确认测试与系统测试 \hfill（第169--171页）}
\begin{itemize}
    \item 确认测试（Validation Testing）—— 验证是否满足需求
    \item 功能测试
    \item \textbf{非功能性测试：}
    \begin{itemize}
        \item 性能测试（Performance Testing）
        \item 压力测试（Stress Testing）
        \item 容量测试（Volume Testing）
        \item 安全性测试（Security Testing）
        \item 恢复测试（Recovery Testing）
        \item 用户界面测试（UI Testing）
        \item 兼容性测试（Compatibility Testing）
        \item 可用性测试（Usability Testing）
    \end{itemize}
\end{itemize}

\subsubsection{7.13\quad 验收测试与回归测试 \hfill（第171--172页）}
\begin{itemize}
    \item 验收测试（Acceptance Testing）
    \item $\alpha$测试（Alpha Testing）—— 开发方内部测试
    \item $\beta$测试（Beta Testing）—— 用户方测试
    \item \textbf{回归测试（Regression Testing）}—— 修改后重新测试
    \begin{itemize}
        \item 目的：确保修改未引入新缺陷
        \item 方法：重新执行原有测试用例
    \end{itemize}
    \item 测试过程与测试文档
\end{itemize}

\subsection{Part 5\quad 软件测试管理 \hfill（第172--175页）}

\subsubsection{7.14\quad 测试组织与计划 \hfill（第172--173页）}
\begin{itemize}
    \item 成立软件测试组织（独立于开发团队）
    \item 软件测试的输出
    \item 制定和实施软件测试计划：
    \begin{itemize}
        \item 标识符、测试项、测试策略
        \item 交付物、任务、环境、进度、风险
    \end{itemize}
\end{itemize}

\subsubsection{7.15\quad 软件测试八项流程 \hfill（第173--175页）}
\begin{enumerate}
    \item 熟悉项目情况
    \item 排序测试需求 —— 测试需求排序示例
    \item 定义测试策略
    \item 设计测试用例
    \item 估计测试工作量
    \item 配置测试资源
    \item 执行测试
    \item 书写测试报告
\end{enumerate}

\vspace{0.3cm}
% ============================================
% 第8章
% ============================================
\section{第8章：软件部署、维护与演化 \hfill（第175--186页）}

\subsection{Part 1\quad 软件部署 \hfill（第175--180页）}

\subsubsection{8.1\quad 软件部署的概念 \hfill（第176页）}
\begin{itemize}
    \item 何为软件部署（收集、打包、安装、配置、发布）
    \item 软件运行环境
    \item 安装和配置运行环境
    \item 安装和配置软件系统
    \item 软件与其运行环境的关系
\end{itemize}

\subsubsection{8.2\quad 部署的挑战与原则 \hfill（第177页）}
\begin{itemize}
    \item 软件间的关系：纵向层次性 / 横向相关性
    \item 软件运行环境的变化
    \item 软件部署的主要问题和挑战
    \item 软件部署的原则：
    \begin{enumerate}
        \item 最小化原则
        \item 相关性原则
        \item 适应性原则
    \end{enumerate}
\end{itemize}

\subsubsection{8.3\quad 部署方式与方法 \hfill（第178--180页）}
\begin{itemize}
    \item 单机部署方式
    \item 分布式部署
    \item 软件部署方法：
    \begin{itemize}
        \item 基于操作系统的部署
        \item 基于软件中间件的部署
        \item 基于容器的部署（Docker）
    \end{itemize}
    \item \textbf{云原生（Cloud Native）}
    \item 基于部署图（Deployment Diagram）规划部署
    \item 部署工具：
    \begin{itemize}
        \item Fat jar（基于语言的部署）
        \item Docker（基于容器的部署）
        \item Jenkins（基于流水线的部署）
    \end{itemize}
\end{itemize}

\subsection{Part 2\quad 软件维护 \hfill（第180--184页）}

\subsubsection{8.4\quad 软件维护的概念与特点 \hfill（第180--181页）}
\begin{itemize}
    \item 软件也需要进行维护
    \item 何为软件维护
    \item 软件维护的特点：同步性、周期长、费用高、难度大
    \item 软件维护成本模型：$M = P + K \cdot e^{(c-d)}$
\end{itemize}

\subsubsection{8.5\quad 软件维护分类 \hfill（第181--182页）}
\begin{enumerate}
    \item \textbf{纠错性维护（Corrective Maintenance）}—— 修复发现的缺陷
    \item \textbf{适应性维护（Adaptive Maintenance）}—— 适应外部环境变化
    \item \textbf{完善性维护（Perfective Maintenance）}—— 提升已有功能和性能
    \item \textbf{预防性维护（Preventive Maintenance）}—— 预防潜在问题
\end{enumerate}

\subsubsection{8.6\quad 维护的问题与可维护性 \hfill（第182--183页）}
\begin{itemize}
    \item 软件维护需要解决的问题：
    \begin{itemize}
        \item 人员的问题
        \item 软件制品的问题
        \item 维护副作用的问题
    \end{itemize}
    \item 软件的可维护性（4个特性）：
    \begin{enumerate}
        \item 可理解性（Understandability）
        \item 可测试性（Testability）
        \item 可修改性（Modifiability）
        \item 可移植性（Portability）
    \end{enumerate}
\end{itemize}

\subsubsection{8.7\quad 软件维护过程 \hfill（第183--184页）}
\begin{itemize}
    \item 软件维护过程
    \item 软件维护工作流程（12步）
    \item 维护管理文档
    \item 维护文档
\end{itemize}

\subsection{Part 3\quad 软件演化 \hfill（第184--186页）}

\subsubsection{8.8\quad 软件逻辑老化 \hfill（第184页）}
\begin{itemize}
    \item 何为软件逻辑老化（Software Aging）
    \item 软件逻辑老化的现象
    \item 软件逻辑老化的原因：
    \begin{itemize}
        \item 缺乏变更
        \item 质量下降
        \item 变更成本增加
        \item 负面变更
        \item 用户满意度降低
    \end{itemize}
\end{itemize}

\subsubsection{8.9\quad 软件演化及解决方法 \hfill（第185--186页）}
\begin{itemize}
    \item 软件演化及其特点（4个特性）
    \item 解决方法：
    \begin{enumerate}
        \item \textbf{代码重组（Code Restructuring）}
        \item \textbf{设计重构（Design Refactoring）}
        \item \textbf{逆向工程（Reverse Engineering）}—— 从代码反推设计
        \item \textbf{再工程（Re-engineering）}—— 重新设计再实现
    \end{enumerate}
    \item 小结
\end{itemize}

\vspace{0.5cm}
\begin{center}
    \rule{\textwidth}{1pt}\\[0.3cm]
    {\large\bfseries —— 全册目录结束 ——}\\[0.3cm]
    {\normalsize 共8章 \quad 186页 \quad 祝考试顺利！}
\end{center}

\end{document}
